\documentclass[10pt]{article}

\usepackage[a4paper,margin=1.9cm]{geometry}
\usepackage{amsmath,amssymb,amsthm}
\usepackage{array}
\usepackage{longtable}
\usepackage{booktabs}
\usepackage{xcolor}
\usepackage{colortbl}
\usepackage[hidelinks]{hyperref}

% ---- notation ---------------------------------------------------------------
\newcommand{\N}{\mathbb{N}}
\newcommand{\Col}{\operatorname{Col}}
\newcommand{\Syr}{\operatorname{Syr}}
\newcommand{\odd}{\operatorname{odd}}
\newcommand{\Colmin}{\Col_{\min}}
\newcommand{\Syrmin}{\Syr_{\min}}
\newcommand{\aal}{\text{a.a.}}
\newcommand{\code}[1]{\texttt{\small #1}}
\newcommand{\yes}{\cellcolor{green!12}}
\newcommand{\hyp}{\cellcolor{red!12}}
\newcommand{\str}{\cellcolor{blue!8}}

% column helpers
\newcolumntype{L}[1]{>{\raggedright\arraybackslash}p{#1}}
\newcolumntype{C}[1]{>{\centering\arraybackslash}p{#1}}

\title{\bfseries A Matrix-Only, Self-Contained Proof of\\
Tao's Theorem 1.3\\[2pt]
\large ``Almost all Collatz orbits attain almost bounded values''}
\author{}
\date{}

\begin{document}
\maketitle
\vspace{-2.2em}

\begin{center}
\emph{Every step of the argument is presented as a matrix.}\\[2pt]
\footnotesize
Legend:\quad
\fbox{\yes\ proved} \quad
\fbox{\str\ structural reduction (proved)} \quad
\fbox{\hyp\ irreducible analytic input}
\end{center}

%==============================================================================
\section*{Matrix 0 \; --- \; Objects and definitions (self-contained)}
%==============================================================================

\renewcommand{\arraystretch}{1.35}
\begin{longtable}{L{2.55cm} L{7.2cm} L{5.8cm}}
\toprule
\textbf{Object} & \textbf{Definition} & \textbf{Notes / Idris} \\
\midrule
\endhead
Collatz map $\Col$ &
$\Col(n)=\begin{cases} n/2 & n \text{ even}\\ 3n+1 & n \text{ odd}\end{cases}$ &
on $\N^{+}$; \code{Dynamics.Col} \\
Odd part $\odd$ &
$\odd(n)=$ the odd number obtained by deleting all factors of $2$ from $n$;
$\odd(n)=n$ if $n$ odd &
$\odd(2^k q)=q$; \code{Dynamics.oddPart}, \code{oddFactor} \\
Syracuse map $\Syr$ &
$\Syr(q)=\odd(3q+1)$, acting on odd $q$ &
$\Syr=\odd\circ\Col$ on odds; \code{Dynamics.Syr} \\
Iteration &
$\Col^{0}=\mathrm{id}$, $\Col^{k+1}=\Col\circ\Col^{k}$ (same for $\Syr$) &
\code{Core.iter}, \code{Core.iterPlus} \\
Orbit minimum &
$\Colmin(n)=\inf_{k\ge 0}\Col^{k}(n)$,\;
$\Syrmin(q)=\inf_{k\ge 0}\Syr^{k}(q)$ &
realised as ``eventually below'' \\
Eventually below &
$\Col{\downarrow}(n,b):\Leftrightarrow \exists k.\ \Col^{k}(n)\le b$ &
\code{Core.EventuallyBelow}, \code{ColBelow}, \code{SyrBelow} \\
$f\to\infty$ &
$\forall t\,\exists T\,\forall x.\ x\ge T \Rightarrow f(x)\ge t$ &
\code{Large.TendsToInfinityOn} \\
Count &
$\#\{\,m<N : p(m)\,\}$, written $\mathrm{cnt}_p(N)$ &
\code{Density.count} \\
Negligible (density $0$) &
$p$ negligible $:\Leftrightarrow \forall k\,\exists N_0\,\forall N\ge N_0.\
(k{+}1)\,\mathrm{cnt}_p(N)\le N$ &
\code{Density.Negligible} \\
Almost all &
$\aal\,p :\Leftrightarrow (\lnot p)$ is negligible &
\code{Density.AlmostAll}; \code{Large.AlmostAllOn} \\
\bottomrule
\end{longtable}

\noindent
The three headline statements are matrices of quantifiers:
\[
\renewcommand{\arraystretch}{1.5}
\begin{array}{r@{\;}c@{\;}l}
\textbf{Thm 1.3 (Collatz)} & : & \forall f\,(f\to\infty)\ \Rightarrow\
   \aal\ n:\ \Col{\downarrow}\big(n,\,f(n)\big),\\[1pt]
\textbf{Thm 1.6 (Syracuse)} & : & \forall g\,(g\to\infty)\ \Rightarrow\
   \aal\ q:\ \Syr{\downarrow}\big(q,\,g(q)\big),\\[1pt]
\textbf{Thm 3.1 (quant.)} & : & \forall g\,(g\to\infty)\ \Rightarrow\
   \{q:\ \lnot\,\Syr{\downarrow}(q,g(q))\}\ \text{has density }0.
\end{array}
\]
(\code{Large.Theorem13}, \code{Large.Theorem16}, \code{Dependencies.QuantitativeSyracuseBound}.)

%==============================================================================
\section*{Matrix 1 \; --- \; The reduction chain (rows compose top to bottom)}
%==============================================================================

\renewcommand{\arraystretch}{1.45}
\begin{longtable}{C{0.7cm} L{3.0cm} L{3.0cm} L{4.3cm} L{3.2cm}}
\toprule
\# & \textbf{From} & \textbf{To} & \textbf{Justification (key idea)} & \textbf{Idris} \\
\midrule
\endhead
\str R1 & Props.\ 1.9 \& 7.8 & Prop.\ 1.11 (stabilisation) &
 combine valuation tail estimate with renewal stability &
 \code{analyticInput\-ToStabilisation} \\
\str R2 & Prop.\ 1.11 & Thm 3.1 &
 iterate the stabilised first-passage bound &
 \code{proposition11To\-Theorem31} \\
\str R3 & Thm 3.1 & Thm 1.6 &
 a density-$0$ exceptional set \emph{is} the ``almost all'' conclusion &
 \code{theorem31To\-Theorem16} \\
\str R4 & Thm 1.6 & Thm 1.3 &
 odd-part orbit simulation $+$ threshold compatibility (Matrix 2--3) &
 \code{theorem16To\-Theorem13Constructive} \\
\midrule
\yes $\Sigma$ & Props.\ 1.9 \& 7.8 & \textbf{Thm 1.3} &
 $\text{R4}\circ\text{R3}\circ\text{R2}\circ\text{R1}$ &
 \code{structuredCentral\-Theorem} \\
\bottomrule
\end{longtable}

\noindent
The composite is exactly the product of the four one-step reductions:
\[
\underbrace{\text{Thm 1.3}}_{\text{Collatz}}
\;\xleftarrow{\ \text{R4}\ }\;
\underbrace{\text{Thm 1.6}}_{\text{Syracuse}}
\;\xleftarrow{\ \text{R3}\ }\;
\underbrace{\text{Thm 3.1}}_{\text{quantitative}}
\;\xleftarrow{\ \text{R2}\ }\;
\underbrace{\text{Prop 1.11}}_{\text{stabilisation}}
\;\xleftarrow{\ \text{R1}\ }\;
\underbrace{\text{Props 1.9, 7.8}}_{\text{analytic input}} .
\]

\noindent\textbf{Dependency (adjacency) matrix} $A$, with $A_{ij}=1$ iff row $i$
is used to prove column $j$ (reading ``$j$ depends on $i$''):
\[
\renewcommand{\arraystretch}{1.2}
\begin{array}{c|ccccc}
 & \text{1.9/7.8} & \text{1.11} & \text{3.1} & \text{1.6} & \text{1.3}\\
\hline
\text{1.9/7.8} & 0 & 1 & 0 & 0 & 0\\
\text{1.11}    & 0 & 0 & 1 & 0 & 0\\
\text{3.1}     & 0 & 0 & 0 & 1 & 0\\
\text{1.6}     & 0 & 0 & 0 & 0 & 1\\
\text{1.3}     & 0 & 0 & 0 & 0 & 0
\end{array}
\qquad
\text{(strictly upper triangular }\Rightarrow\text{ acyclic; a single directed path).}
\]

%==============================================================================
\section*{Matrix 2 \; --- \; Step R4 proved: the odd-part orbit simulation}
%==============================================================================

\noindent
R4 is the only step with genuine dynamical content; it is fully proved.  Its
skeleton is the matrix of facts below (all in \code{OddPart}/\code{OddThreshold}).

\renewcommand{\arraystretch}{1.4}
\begin{longtable}{C{0.7cm} L{7.6cm} L{5.7cm}}
\toprule
\# & \textbf{Fact} & \textbf{Idris} \\
\midrule
\endhead
\yes a & $n$ odd $\Rightarrow \Col(n)=3n+1$; $n$ even $\Rightarrow \Col(n)=n/2$ &
 \code{colOddStep}, \code{colEvenStep} \\
\yes b & Deleting factors of $2$ terminates: $\odd(n)$ is odd for $n\ge 1$ &
 \code{oddFactorOdd}, \code{syrIsOdd} \\
\yes c & Iterating $\Col$ from $n$ reaches $\odd(n)$ in $\code{oddPartDropTime}(n)$
 steps, with all intermediate heights $\ge \odd(n)$ &
 \code{oddPartDropTime\-NormalizesValue} \\
\yes d & One $\Syr$-step is realised by finitely many $\Col$-steps:
 $\exists s.\ \Col^{s}(m)=\Syr^{t}(\odd(m))$ (as values) &
 \code{syrRealize}, \code{syrRealizeStep} \\
\yes e & Hence a height-dominating semiconjugacy
 $\odd:(\N^{+},\Col)\to(\text{odds},\Syr)$: for all $n,t$
 $\exists s.\ \Col^{s}(n)\le \Syr^{t}(\odd(n))$ &
 \code{provenOddPart\-OrbitSimulation} \\
\yes f & Simulations form a category (identity, composition) &
 \code{orbitSimulationId}, \code{orbitSimulationCompose} \\
\bottomrule
\end{longtable}

\noindent
\textbf{Consequence (transfer).}\;
If $\Syr{\downarrow}(\odd(n),b)$ then $\Col{\downarrow}(n,b)$: pick $t$ with
$\Syr^{t}(\odd(n))\le b$; by row (e) some $\Col^{s}(n)\le \Syr^{t}(\odd(n))\le b$.
\hfill\code{oddPartOrbitTransfer}

%==============================================================================
\section*{Matrix 3 \; --- \; Step R4 proved: the threshold is \emph{constructed}}
%==============================================================================

\noindent
To run Thm 1.6 with a Syracuse growth function whose value at $\odd(n)$ never
exceeds $f(n)$, the threshold is \emph{built} from the growth witness
$w$ of $f$ (no separate hypothesis).  Define, for odd $q$,
\[
   t_w(q)\;=\;\max\{\,t\le q :\ \code{thresholdFor}(w)(t)\le q\,\}.
\]

\renewcommand{\arraystretch}{1.4}
\begin{longtable}{C{0.7cm} L{8.6cm} L{4.7cm}}
\toprule
\# & \textbf{Property} & \textbf{Idris} \\
\midrule
\endhead
\yes g & $\odd(n)\le n$ (deleting factors of $2$ cannot increase) &
 \code{oddFactorLe}, \code{oddPartSizeLe} \\
\yes h & Compatibility: $t_w(\odd(n))\le f(n)$ &
 \code{oddThresholdOfCompatible} \\
\yes i & Growth: $f\to\infty \Rightarrow t_w\to\infty$ &
 \code{oddThresholdOfGrows} \\
\yes j & Bounded search is correct (largest admissible target) &
 \code{findBest}, \code{findBestCompat}, \code{findBestGe} \\
\bottomrule
\end{longtable}

\noindent
\textbf{Proof of R4.}\;
Given $g:=t_w$, Thm 1.6 yields $\aal\ q:\ \Syr{\downarrow}(q,t_w(q))$.  Pull
back along $\odd$ (Matrix 4, closure) and monotonicity in the bound (row h) gives
$\aal\ n:\ \Syr{\downarrow}(\odd(n),f(n))$; Matrix 2 (transfer) converts each to
$\Col{\downarrow}(n,f(n))$.  \hfill$\square$\;\;\code{liftSyracuseToCollatz}

%==============================================================================
\section*{Matrix 4 \; --- \; The ``almost all'' algebra (density model, proved)}
%==============================================================================

\noindent
``Almost all'' means the exceptional set has natural density $0$.  This is a
genuine, non-degenerate measure of smallness; the closure laws it must obey are
all proved from first principles.

\renewcommand{\arraystretch}{1.4}
\begin{longtable}{L{5.3cm} L{6.2cm} L{3.5cm}}
\toprule
\textbf{Law} & \textbf{Statement} & \textbf{Idris} \\
\midrule
\endhead
Monotone (subset) & $S\subseteq T$, $S$ negligible only via $T$; a.a.\ closed under supersets &
 \code{negligibleMono}, \code{almostAllMono} \\
Finite $\cup$ / $\cap$ & union of negligibles is negligible; intersection of a.a.\ is a.a. &
 \code{orNegligible}, \code{andAlmostAll}, \code{orListNegligible}, \code{andListAlmostAll} \\
Bounded $\Rightarrow$ small & finite / cofinite sets are negligible / a.a. &
 \code{boundedNegligible}, \code{almostAllCofinite} \\
Complementary count & $\mathrm{cnt}_p(N)+\mathrm{cnt}_{\lnot p}(N)=N$ &
 \code{countComplement} \\
Full space & $\mathrm{cnt}_{\top}(N)=N$; \; density$(\top)=1$ &
 \code{countAllTrue} \\
\rowcolor{green!12} Non-degeneracy & the whole space is \textbf{not} negligible (else $2N\le N$) &
 \code{allNotNegligible} \\
Cofinal / teeth & a negligible set omits arbitrarily large $n$ &
 \code{negligibleCofalse}, \code{almostAllCofinal} \\
On the carriers & same algebra on $\N^{+}$ (\code{Pos}) and odds (\code{OddPos}) &
 \code{CarrierDensity.*} \\
\bottomrule
\end{longtable}

\noindent
Pull-back and push-forward used by R3, R4 are the structural maps
\code{controlPullback}, \code{controlMap} (\code{Large}); combined with the laws
above they move ``almost all'' along $\odd$ and along pointwise implications.

%==============================================================================
\section*{Matrix 5 \; --- \; What is proved vs.\ what is assumed}
%==============================================================================

\renewcommand{\arraystretch}{1.4}
\begin{longtable}{L{4.6cm} C{2.2cm} L{7.8cm}}
\toprule
\textbf{Ingredient} & \textbf{Status} & \textbf{Remark} \\
\midrule
\endhead
Collatz/Syracuse dynamics & \yes proved & computable; $\Col(3)=10$, $\Syr(7)=11$ \\
Odd-part correspondence (R4) & \yes proved & Matrices 2--3, no \code{believe\_me} \\
Threshold system & \yes proved & \emph{constructed} from the growth witness \\
Reduction plumbing R1--R3 & \str proved & structural; transports its input verbatim \\
``Almost all'' algebra & \yes proved & density model, non-degenerate (Matrix 4) \\
\rowcolor{red!12} Prop.\ 1.9 (valuation tail) & assumed & $2$-adic valuation distribution \& sub-Gaussian tail \\
\rowcolor{red!12} Prop.\ 7.8 (renewal stability) & assumed & monotonicity of the renewal process \\
\midrule
\rowcolor{green!12}\textbf{Thm 1.3} & \textbf{proved from the two inputs} &
 \code{structuredCentralTheorem}; provably equal to \code{centralTheoremUnconditional} \\
\bottomrule
\end{longtable}

\begin{center}\footnotesize
The single remaining mathematical dependency is the pair of deep analytic
estimates (Props.\ 1.9 and 7.8), highlighted in red.  Everything else --- the
entire logical skeleton, the odd-part dynamics, the constructed threshold, and
the non-degenerate density model of ``almost all'' --- is a total, machine-checked
Idris proof containing no \code{believe\_me} and no axioms.
\end{center}

\end{document}
